\section{Limitations and Conclusion}

In this paper, we present a methodology for creating instruction-following foundation models of behavior. Specifically, by leveraging two existing pretrained foundation models: a behavioral prior (VPT \citep{baker2022vpt}) and a domain-specific CLIP model (MineCLIP \citep{fan2022minedojo}), we create a powerful Minecraft agent that can follow short-horizon open-ended text and visual instructions, all for only \$60 of compute. The resulting foundation model, \agentname{}, sets a new bar for open-ended instruction-following in Minecraft with low-level controls (mouse and keyboard) and raw pixel inputs, far outperforming previous baselines and robustly completing 12 of 13 tasks in our early-game evaluation suite. We note that generalist agents such as \agentname{} can have potential negative effects on society. We include a thorough discussion of these issues in \Cref{appdx:broader_impact}.

\agentname{} is a significant advancement in creating generative models of text-to-behavior, but it has several limitations, as described in \Cref{appx:limitations_future}. First, \agentname{} is mostly proficient at achieving short-horizon tasks while struggling with longer-horizon tasks. While prompt chaining is a promising approach for improving performance on complex tasks, more can be done in future work to improve performance. Another limitation we observe is that prompt engineering, as with other generative models, can be unintuitive and time-consuming. Future work should investigate improving the steerability of \agentname{} through a better understanding of natural language prompts. Additionally, we note that evaluating and describing the capabilities of open-ended generalist agents is an open research problem itself since capability depends strongly on preconditions, prompt engineering, and our own ability to come up with varied and challenging tasks. Finally, since our approach is not specific to the Minecraft domain, we hope that the method used to create \agentname{} can inspire future work in creating powerful generalist agents in other domains and environments.
